\begin{table}[p]
\centering
\begin{minipage}[t]{.55\textwidth}
\centering
\begin{tabular}{rrrrr}
\toprule
$B<P$ & $B<Q$ & $P<H$ & $Q<H$ & Count\\
\midrule
2T & 2T & 2T & 2T & 4\\
2T & 2T & M & M & 21\\
2T & 2T & X & X & 315\\
2T & M & M & 2T & 1\\
2T & M & X & X & 32\\
M & 2T & 2T & M & 1\\
M & 2T & X & X & 32\\
M & M & 2T & 2T & 1\\
M & M & M & M & 88\\
M & M & X & X & 587\\
\bottomrule
\end{tabular}
\end{minipage}\hfill
\begin{minipage}[t]{.41\textwidth}
\centering
\begin{tabular}{rrrrr}
\toprule
$H$ & $P$ & $Q$ & $B$ & Count\\
\midrule
N & N & N & Y & 1\\
Y & N & N & N & 404\\
Y & N & N & Y & 674\\
Y & N & Y & N & 1\\
Y & Y & N & N & 1\\
Y & Y & Y & Y & 1\\
\bottomrule
\end{tabular}
\end{minipage}
\caption{Cubic symmetric Max and Norm counts.  Every displayed count equals the corresponding printed entry of \cite{L059}.  The cubic table uses the bipartite symmetric graphs through 2048 vertices.}
\end{table}
\begin{table}[p]
\centering
\begin{minipage}[t]{.55\textwidth}
\centering
\begin{tabular}{rrrrr}
\toprule
$B<P$ & $B<Q$ & $P<H$ & $Q<H$ & Count\\
\midrule
2T & 2T & 2T & 2T & 10\\
2T & 2T & M & M & 5\\
2T & 2T & X & X & 38\\
2T & X & M & M & 1\\
2T & X & X & 2T & 11\\
2T & X & X & M & 1\\
2T & X & X & X & 19\\
X & 2T & 2T & X & 11\\
X & 2T & X & X & 26\\
X & X & 2T & 2T & 1\\
X & X & M & M & 17\\
X & X & X & X & 1503\\
\bottomrule
\end{tabular}
\end{minipage}\hfill
\begin{minipage}[t]{.41\textwidth}
\centering
\begin{tabular}{rrrrr}
\toprule
$H$ & $P$ & $Q$ & $B$ & Count\\
\midrule
N & N & N & N & 365\\
N & N & N & Y & 268\\
N & N & Y & N & 4\\
N & N & Y & Y & 3\\
N & Y & N & N & 4\\
N & Y & N & Y & 3\\
Y & N & N & N & 466\\
Y & N & N & Y & 525\\
Y & N & Y & N & 2\\
Y & Y & N & N & 2\\
Y & Y & Y & Y & 1\\
\bottomrule
\end{tabular}
\end{minipage}
\caption{Quartic Max and Norm counts.  Every displayed count equals the corresponding printed entry of \cite{L059}.  The quartic threshold and ordered-part correction are explicit in the text.}
\end{table}
